% titel: Blatt 6, A5: Welche Verbindungen sind nicht bildbar?
% skript: Mortimer Kap. 8 (chemie: Atombindung, Periodengrenzen)

Ob eine Formel ein denkbares Molekül ergibt, hängt an der maximalen
Bindungszahl der Zentralatome. Elemente der 2.\ Periode können ihr Oktett
niemals erweitern; C bildet höchstens 4, N höchstens 3 (bzw.\ mit
Donor-Koordination 4) Bindungen.

\begin{werkzeugkasten}[Kriterien für bildbare Verbindungen]
\begin{itemize}
  \item Zentralatom der 2.\ Periode: keine $d$-Orbitale $\to$ keine
        Oktetterweiterung (N, C, O können nicht $>8$ Valenzelektronen haben).
  \item C: maximal 4 gültige Bindungen (Vier-Bindungsgrenze).
  \item Zentralatome ab Periode~3 (P, S, Se, I) \textit{dürfen} ihr Oktett
        erweitern $\to$ mehr als 4 Bindungen möglich.
  \item Elektronegative Endgruppen mit grosser EN-Differenz stabilisieren
        hohe Koordinationszahlen (z.\ B. $\mathrm{Cl}$).
\end{itemize}
\textit{Skript:} Mortimer Kap. 8 (Atombindung, Oktettregel).
\end{werkzeugkasten}

\begin{achtung}[Darauf achten]
\begin{itemize}
  \item Prüfen, in welcher Periode das \textit{Zentralatom} steht --- danach
        richtet sich, ob Erweiterung möglich ist.
  \item $\mathrm{CH_5}$ würde vierbindiges C mit zusätzlichem H beanspruchen;
        C kann das nicht (und Elektronen zählen: nur 8).
  \item Der Gegensatz ist der klassische Stolperstein:
        $\mathrm{NCl_3}$ ja, aber $\mathrm{NCl_5}$ nein --- N ist keine
        Gruppe-15-Erweiterung fähig wie P.
\end{itemize}
\end{achtung}
